\chapter[Characters and Schur polynomials]{Characters and Schur Polynomials}

Cycle lengths describe the fixed tabloids of a permutation. An
alternating coefficient extraction then separates the irreducible
characters. The same coefficients relate these characters to symmetric
polynomials and turn induction into multiplication.

\section{Cycle counting and the Frobenius formula}

For \(\nu\vdash n\), write \(m_j(\nu)\) for the multiplicity of \(j\).
Use \(z_\nu=\prod_j j^{m_j(\nu)}m_j(\nu)!\), as in
Result~\ref{res:A01}. In \(r\) variables set
\[
 p_j(x)=\sum_{i=1}^r x_i^j,\qquad
 p_\nu(x)=\prod_{j\ge1}p_j(x)^{m_j(\nu)}.
\]
All products over an empty set are \(1\). The variables in coefficient
extractions commute, and an exponent with a negative component has
coefficient zero.

\begin{proposition}[The Young permutation character]\label{res:C01}
Let \(\mu=(\mu_1,\ldots,\mu_r)\) be a list of nonnegative integers summing
to \(n\). Let \(\xi^\mu\) be the character of the \(S_n\)-module with
basis the ordered decompositions of the labels into rows of these sizes.
For cycle type \(\nu\), one has \(\xi^\mu(\nu)=[x^\mu]p_\nu(x)\).
\end{proposition}

\begin{proof}
This is the tabloid module of Result~\ref{res:A02}, with ordered row
sizes. Reordering the sizes gives an equivariant bijection by
reordering the rows. The trace on the basis counts fixed decompositions.
A row is preserved by the permutation precisely when it is a union of
cycles. A cycle of length \(j\) can therefore be assigned to any row
\(i\), contributing \(x_i^j\). The product records independent assignments
of the individual cycles, and coefficient extraction imposes the sizes
of all rows. See \cite[Section 4.14]{etingof}.
\end{proof}

For \(r\ge1\), let \(\delta=(r-1,r-2,\ldots,0)\). For a vector \(b\)
of \(r\) nonnegative integers define \(a_b(x)=\det[x_i^{b_j}]_{i,j=1}^r\).
Thus \(\Delta(x)=a_\delta(x)=\prod_{i<j}(x_i-x_j)\). The row convention
in this determinant fixes all the signs in the coefficient formula.

\begin{lemma}[Cauchy kernel and alternating coefficients]\label{res:C02}
For \(r\) variables in each alphabet,
\begin{align*}
 \prod_{i,j}(1-x_i y_j)^{-1}
   &=\sum_\nu p_\nu(x)p_\nu(y)/z_\nu,\\
 \Delta(x)\Delta(y)\prod_{i,j}(1-x_i y_j)^{-1}
   &=\sum_{\ell(\nu)\le r}a_{\nu+\delta}(x)a_{\nu+\delta}(y).
\end{align*}
Consequently, the coefficient of \(x^{\lambda+\delta}y^{\mu+\delta}\)
in the second line is \(1\) if \(\lambda=\mu\) and \(0\) otherwise.
\end{lemma}

\begin{proof}
The first identity follows by expanding
\(\exp(\sum_{j\ge1}p_j(x)p_j(y)/j)\): using the \(j\)-th summand \(m\)
times contributes \(1/(j^m m!)\). All operations take place degree by
degree in formal power series.

For the second identity, clear the product of denominators in the
determinant of the matrix \([(1-x_i y_j)^{-1}]\).
The resulting polynomial alternates separately in \(x\) and in \(y\),
and its degree in each individual \(x_i\) or \(y_j\) is at most \(r-1\).
It is thus a constant times \(\Delta(x)\Delta(y)\). To determine the
constant, set \(x_r=y_r=0\) and subtract the last row from each preceding
row in the determinant before clearing denominators. The last column is
now zero except for its last entry \(1\), and the remaining entries are
\(x_i y_j/(1-x_i y_j)\). Factoring out these row and column factors
reduces the constant to the one for \(r-1\), with value \(1\) for \(r=1\).

Expand each entry as \(\sum_{k\ge0}x_i^k y_j^k\). For a finite truncation,
the rectangular factors have entries \(x_i^k\) and \(y_j^k\) in the
complex polynomial ring. Apply the rectangular Cauchy--Binet identity
over its field of fractions. Both sides are polynomials, so the
identity holds in the polynomial ring by injectivity of the embedding.
This gives a sum of products of determinants indexed by distinct
exponents \(b_1>\cdots>b_r\ge0\).
Using decreasing order in both determinants preserves their product.
In any fixed degree the sum stabilizes. Such exponent lists are
exactly \(\nu+\delta\) for the stated partitions. A monomial with
strictly decreasing exponents can occur only in the corresponding
determinant and then has coefficient \(1\). This proves the extraction
assertion. Compare \cite[Section 4.15]{etingof}.
\end{proof}

\begin{theorem}[Frobenius character formula]\label{res:C03}
For \(\lambda,\nu\vdash n\), put \(\chi^\lambda=\chi_{S^\lambda}\).
If \(r\ge\ell(\lambda)\) and \(r\ge1\), then
\[
 \chi^\lambda(\nu)
   =[x^{\lambda+\delta}]\Delta(x)p_\nu(x),
\]
where \(\lambda\) is padded with zeros to length \(r\). For the empty
partition the formula also has value \(1\) with zero variables.
\end{theorem}

\begin{proof}
First choose \(r\ge \max(1,n)\). As a function of \(\nu\), the indicated
coefficient is \(\xi^\lambda\) plus an integer combination of
\(\xi^\mu\) for shapes \(\mu\) strictly dominating \(\lambda\).
Indeed the determinant expansion takes signed coefficients with
exponents \(\lambda+\delta-w\delta\). When these exponents are
nonnegative, sort them in decreasing order to obtain \(\mu\).
The sum of the largest \(k\) exponents is at least the sum of the
first \(k\) exponents in the unsorted list. The latter is at least
\(\sum_{i=1}^k\lambda_i\), since the first \(k\) entries of \(\delta\)
have maximal sum among its \(k\)-element subsets. If \(w\ne1\), some
initial subset changes, so at least one of the inequalities is strict.
Result~\ref{res:C01} identifies the corresponding coefficient with
\(\xi^\mu\).

By Results~\ref{res:A06} and \ref{res:A07}, \(\xi^\lambda\) contains
\(\chi^\lambda\) once and only strictly dominating shapes can also
occur. Thus our candidate list is unitriangular in the irreducible
characters, in any order with dominating shapes first. Its entries are
orthonormal: the inner product of two candidates is extraction at
\((\lambda+\delta,\mu+\delta)\) in the first line of
Result~\ref{res:C02} multiplied by \(\Delta(x)\Delta(y)\).
Here the candidates are real, and \(1/z_\nu\) is the weight of the
class in the character inner product. The second line gives the desired
orthonormality. Reading the unitriangular list from most dominant to
least, orthogonality to all earlier entries forces all off-diagonal
coefficients to vanish. We used the finite-group character
orthogonality established in \cite[Chapter 8]{algebra-core}.

To remove a final zero part, specialize the final variable to zero.
The power sums are unchanged, and
\(\Delta_{r+1}(x,0)=(x_1\cdots x_r)\Delta_r(x)\).
The desired coefficient in \(r+1\) variables has last exponent \(0\)
and all other exponents raised by \(1\); it therefore reduces to the
stated \(r\)-variable coefficient. See \cite[Theorem 4.47]{etingof}.
\end{proof}

\section{Schur polynomials and the characteristic map}

For \(r\) variables let \(h_k\) be the sum of the weights of weakly
increasing lists of length \(k\) in \(\{1,\ldots,r\}\), where a letter
\(i\) has weight \(x_i\). Set \(h_0=1\) and \(h_k=0\) for \(k<0\).
Write \(e_k\) for the corresponding sum over strictly increasing lists,
and put \(e_0=1\).

\begin{theorem}[Three formulas for the Schur polynomial]\label{res:C04}
For a partition \(\lambda\) with at most \(r\) parts,
\[
 s_\lambda(x)
 =\frac{a_{\lambda+\delta}(x)}{\Delta(x)}
 =\det[h_{\lambda_i-i+j}(x)]_{i,j=1}^r
 =\sum_T \prod_{u\in\lambda}x_{T(u)} .
\]
Here \(T\) ranges over the tableaux with entries in
\(\{1,\ldots,r\}\), weakly increasing along rows and strictly
increasing down columns. For \(\ell(\lambda)>r\) the tableau sum and
the Jacobi--Trudi determinant using at least \(\ell(\lambda)\) rows
are zero. The empty partition gives \(1\).
\end{theorem}

\begin{proof}
The quotient is a symmetric polynomial because the alternating
numerator is divisible by every \(x_i-x_j\). Divide the second identity
of Result~\ref{res:C02} by \(\Delta(x)\). Extract the coefficient of
\(y^{\lambda+\delta}\) in the resulting equality. On the sum side this
gives \(s_\lambda(x)\). On the other side, expand \(\Delta(y)\) and use
\(\prod_{i}(1-x_i y_j)^{-1}=\sum_{k\ge0}h_k(x)y_j^k\) for each \(j\).
The coefficient of the \(i\)-th variable is \(h_{\lambda_i-i+j}\) in
the term using the \(j\)-th exponent of \(\delta\); the alternating sum
is the stated determinant.

For the tableau formula, use \(r\ge1\) and take a finite rectangular
piece of the integer grid containing all starts \(A_i=(-i,1)\) and
ends \(B_j=(\lambda_j-j,r)\). Use east edges of weight \(x_v\) at
height \(v\) and north edges of weight \(1\). The path sum from \(A_i\)
to \(B_j\) is \(h_{\lambda_j-j+i}\): the heights of the east steps are
weakly increasing and determine a unique path. In zero horizontal
distance there is a single north path; negative distance gives none.
The matrix is the transpose of the Jacobi--Trudi matrix. The weighted
path determinant of \cite[Chapter 3]{finite-combinatorics} writes its
determinant as the signed sum over vertex-disjoint families. The grid
is acyclic, and the starts and the ends are each distinct. Both
boundary lines order endpoints from right to left by their index.
An inversion of that order forces intersection: at each height the
vertices of a path form an interval, and the north steps send its
right endpoint to the left endpoint of its next interval. Two disjoint
intervals preserve their order at every height, including the case
of one height. Hence only the identity matching survives.

For identity matching, the east-step heights in path \(i\) record
row \(i\) of a tableau. For adjacent paths let \(u_j,v_j\) be their
respective \(j\)-th heights. Where both rows have column \(j\), at
horizontal coordinate \(j-i-1\) the first path's vertical interval
is \([u_{j-1},u_j]\) and the second path's is \([v_j,v_{j+1}]\).
Use \(u_0=1\) and the final height \(r\) after the last step.
These are exactly the shared horizontal coordinates of adjacent
paths, since \(\lambda_i\ge\lambda_{i+1}\).
For \(j=1\) the first interval begins at \(1\); disjointness forces
\(u_1<v_1\). Inductively \(u_{j-1}<v_{j-1}\le v_j\), so at the next
coordinate disjointness again forces \(u_j<v_j\). The converse follows
from these same intervals, and their ordered separation prevents
intersection between nonadjacent paths as well: each intermediate
path is present at every shared horizontal coordinate of the two
outer paths. This gives the claimed weight-preserving bijection.
The argument works with any number of rows in the determinant once
it is at least the number of nonzero parts. A column taller than \(r\)
admits no strict filling. With no variables only \(h_0\) is nonzero;
the empty diagram gives \(1\), and for a nonempty shape the first row
of the determinant is zero. Compare \cite[Sections 4.4--4.6]{sagan}.
\end{proof}

Use the ring of symmetric functions in the following degreewise
sense. A degree \(n\) symmetric function is a compatible sequence of
homogeneous symmetric polynomials under setting the last variable to
zero. Its coefficients are determined by the monomial sums for
partitions of \(n\), which stabilize once there are at least \(n\)
variables. The tableau formula shows that \(s_\lambda\) is compatible
with this specialization. The products \(p_\nu\) are a basis over
\(\mathbb C\):
their monomial expansion first assigns the parts of \(\nu\) to
distinct variables, with nonzero coefficient \(\prod_j m_j(\nu)!\);
all other terms merge parts and have fewer distinct variables.
This gives a triangular change of basis. Define the Hall Hermitian
inner product by \(\langle p_\nu,p_\mu\rangle
=\boldsymbol{1}_{\nu=\mu}z_\nu\), linear in its first variable,
with different degrees orthogonal.

\begin{theorem}[Frobenius characteristic]\label{res:C05}
For a class function \(f\) on \(S_n\), the correspondence
\[
 \operatorname{ch}(f)=\sum_{\nu\vdash n}f(\nu)p_\nu/z_\nu
\]
is a degreewise isometric isomorphism and sends \(\chi^\lambda\) to
\(s_\lambda\). For class functions on \(S_a\) and \(S_b\), the
characteristic of their external tensor product induced from
\(S_a\times S_b\) to \(S_{a+b}\) is the product of their
characteristics.
\end{theorem}

\begin{proof}
The class sizes in Result~\ref{res:A01} and the defining norms of the
power sums give the isometry and bijection. For a partition \(\nu\),
the polynomial \(\Delta p_\nu\) is alternating. Its coefficients on
strictly decreasing exponent lists determine it, since lists with a
repeated exponent have zero coefficient. By Result~\ref{res:C03},
expansion in the \(a_{\lambda+\delta}\) and division by \(\Delta\)
give \(p_\nu=\sum_{\lambda}\chi^\lambda(\nu)s_\lambda\) in enough
variables. Orthogonality, as used in Result~\ref{res:C03}, inverts
this matrix and gives the claimed image of \(\chi^\lambda\).

For multiplication, a coset of \(S_a\times S_b\) is an ordered
decomposition into two label blocks, and the cosets fixed by a
permutation are assignments of whole cycles to the two blocks.
The trace on the summand for such a coset is \(f(\alpha)g(\beta)\),
where \(\alpha,\beta\) are the resulting cycle types. For fixed
types, the number of those assignments is
\(\prod_j\binom{m_j(\nu)}{m_j(\alpha)}=z_\nu/(z_\alpha z_\beta)\).
After division by \(z_\nu\), summing over the types is precisely the
product rule \(p_\alpha p_\beta=p_{\alpha\cup\beta}\). The trace
argument proves this for characters, hence by linearity for all
class functions. See \cite[Sections 4.6--4.7]{sagan}.
\end{proof}

\begin{corollary}[Young's rule]\label{res:C06}
The multiplicity of \(S^\lambda\) in the Young permutation module of
content \(\mu\) is \([x^\mu]s_\lambda(x)\). This is the number of
semistandard tableaux of shape \(\lambda\) and content \(\mu\).
\end{corollary}

\begin{proof}
Take the inner product of \(\xi^\mu\) with \(\chi^\lambda\).
By Result~\ref{res:C01} its value is the coefficient of \(x^\mu\) in
\(\sum_\nu\chi^\lambda(\nu)p_\nu(x)/z_\nu\), since the irreducible
character is real by Result~\ref{res:C03}. This polynomial is
\(s_\lambda\) by Result~\ref{res:C05}. Apply the tableau formula in
Result~\ref{res:C04}. See \cite[Theorem 2.11.2]{sagan}.
\end{proof}

\Needspace{7\baselineskip}
\section{Pieri and border-strip rules}

\begin{theorem}[Pieri rules]\label{res:C07}
The product \(s_\lambda h_k\) is the sum of \(s_\mu\) over
\(\mu\supseteq\lambda\) for which \(\mu/\lambda\) has \(k\) cells
with at most one in each column. For \(s_\lambda e_k\), the
corresponding condition is at most one in each row. Each summand
has coefficient \(1\). By the characteristic map these are the
decompositions obtained by inducing the external tensor product of
\(S^\lambda\) with the trivial or sign representation of \(S_k\),
respectively.
\end{theorem}

\begin{proof}
Work in enough variables for the degree in question and put
\(b=\lambda+\delta\). The horizontal-strip condition says
\(\mu_j\ge\lambda_j\) for every \(j\) and
\(\mu_j\le\lambda_{j-1}\) for \(j\ge2\): if that inequality fails,
column \(\lambda_{j-1}+1\) gains cells in rows \(j-1\) and \(j\).
Conversely, if column \(q\) gains a cell in row \(j\), then
\(q\le\mu_j\le\lambda_{j-1}\), so every earlier cell of that column
was already present. In strict
exponents \(c=\mu+\delta\), the independent ranges are \(c_1\ge b_1\)
and \(b_j\le c_j<b_{j-1}\) for \(j\ge2\). Therefore
\(\sum_c t^{\sum_j(c_j-b_j)}a_c(x)\) is the determinant obtained by
summing these geometric series in each column separately. Extract
\((1-x_i t)^{-1}\) from row \(i\). The new first column is \(x_i^{b_1}\);
column \(j\ge2\) is \(x_i^{b_j}-t^{b_{j-1}-b_j}x_i^{b_{j-1}}\).
The column-change matrix from \([x_i^{b_j}]\) is triangular with
diagonal \(1\). The determinant is thus
\(a_b(x)\prod_i(1-x_it)^{-1}\). Extracting \(t^k\) and dividing by
\(\Delta\) gives the first formula.

For the vertical-strip formula, use the alternating sum of
monomials defining \(a_b\). Multiplication by the symmetric
polynomial \(e_k\) commutes with alternation. Before alternating,
each term increments \(k\) different exponents by \(1\). A repeated
exponent gives zero, and the remaining lists stay strictly decreasing
because adjacent exponents originally differed by at least \(1\).
They give exactly the partitions with one added cell in each of
\(k\) distinct rows. Finally a single row has Schur function \(h_k\)
and a single column has \(e_k\) by Result~\ref{res:C04}. The last
assertion follows from Results~\ref{res:A07} and \ref{res:C05}.
The case \(k=1\) agrees with Result~\ref{res:B05}.
Compare \cite[Section 4.9]{sagan}.
\end{proof}

\begin{theorem}[Murnaghan--Nakayama rule]\label{res:C08}
For \(m\ge1\) and \(\lambda\vdash n\) with \(m\le n\), let
\(\tau\vdash n-m\). Then
\[
 \chi^\lambda((m)\cup\tau)
   =\sum_{\substack{\mu\subseteq\lambda\\
             \lambda/\mu\ \text{a border strip of size }m}}
        (-1)^{\operatorname{ht}(\lambda/\mu)}\chi^\mu(\tau).
\]
A border strip is connected by edges and contains no \(2\)-by-\(2\)
square. Its height is its number of occupied rows minus one.
The empty character takes value \(1\) on the empty cycle type.
\end{theorem}

\begin{proof}
For a partition \(\mu\), set \(b=\mu+\delta\). Multiplying the
alternating monomial sum for \(a_b\) by \(p_m\) adds \(m\) to one
exponent \(b_j\) before alternation. Discard coincident exponents.
Move the raised exponent to its place \(i\le j\) in the strictly
decreasing list \(c\); each interchange contributes a sign, giving
overall \((-1)^{j-i}\). The new partition \(\lambda=c-\delta\)
agrees with \(\mu\) outside rows \(i,\ldots,j\), increases in row
\(i\), and has
\(\lambda_k=\mu_{k-1}+1\) for \(i<k\le j\). All these rows gain
cells, and consecutive added row intervals meet in exactly one
column. Thus the skew shape is a border strip of height \(j-i\).
Conversely, for a border strip with occupied rows \(i,\ldots,j\),
connectedness gives an overlap at every adjacent pair of rows.
Because there is no square that overlap has length one.
The added intervals start at \(\mu_k+1\), so this forces
\(\lambda_k=\mu_{k-1}+1\) for \(i<k\le j\).
The total of \(m\) added cells determines \(\lambda_i\), recovering
exactly the exponent move above. Hence
\(p_m s_\mu=\sum_\lambda(-1)^{\operatorname{ht}(\lambda/\mu)}s_\lambda\).

The coefficient of \(s_\lambda\) in \(p_{(m)\cup\tau}=p_m p_\tau\)
is \(\chi^\lambda((m)\cup\tau)\) by the expansion in
Result~\ref{res:C05}. Expand \(p_\tau\) first and apply the
edge-strip multiplication just proved. This gives the rule,
including the initial value. See \cite[Section 4.10]{sagan}.
\end{proof}

\paragraph{Example: the character table of \(S_4\).}
Class sizes follow from Result~\ref{res:A01}. In the standard
representation the trace is the number of fixed points minus one
by Result~\ref{res:B09}. Its sign twist gives the row \((2,1,1)\).
For \((2,2)\), a \(2\)-strip leaves either a horizontal pair with
positive sign or a vertical pair with negative sign. A \(3\)-strip
leaves one cell with negative sign, and the full square is no
\(4\)-strip. The value at the identity is \(2\) by
Result~\ref{res:B08}. These computations give
\[
\begin{array}{c|rrrrr}
 \nu & 1^4 & (2,1,1) & (2,2) & (3,1) & (4)\\ \hline
 \text{class size} & 1 & 6 & 3 & 8 & 6\\ \hline
 (4)     &1& 1& 1& 1& 1\\
 (3,1)   &3& 1&-1& 0&-1\\
 (2,2)   &2& 0& 2&-1& 0\\
 (2,1,1) &3&-1&-1& 0& 1\\
 (1^4)   &1&-1& 1& 1&-1
\end{array}
\]
With the indicated class-size weights the row Gram matrix is
\(24I\). The squared column norms are \(24,4,8,3,4\), the five
centralizer orders, and distinct columns are orthogonal.

\Needspace{26\baselineskip}
\section{The dual Jacobi--Trudi formula}

\begin{corollary}[Dual Jacobi--Trudi]\label{res:C09}
Let \(\lambda^{\mathsf t}\) be the transpose of the partition
\(\lambda\), and pad it with zeros to length \(q\ge \lambda_1\)
with \(\lambda_1=0\) for the empty partition.
Then in any finite alphabet \(x\),
\[
 s_\lambda(x)=
 \det\bigl[e_{\lambda^{\mathsf t}_i-i+j}(x)\bigr]_{i,j=1}^{q}.
\]
Here \(e_0=1\) and \(e_k=0\) for \(k<0\). For the empty shape one
may take \(q=0\) and use the empty determinant \(1\).
\end{corollary}

\begin{proof}
On the stable symmetric functions used in \ref{res:C05}, define
\(\omega(p_j)=(-1)^{j-1}p_j\), extending multiplicatively and
linearly. The products \(p_\nu\) are a basis, so this rule defines
an algebra involution. Both generating series are identities
in stable symmetric functions:
\begin{align*}
 \sum_{k\ge0}h_k t^k &= \exp\!\left(\sum_{j\ge1}p_jt^j/j\right),\\
 \sum_{k\ge0}e_k t^k &= \exp\!\left(\sum_{j\ge1}(-1)^{j-1}p_jt^j/j\right).
\end{align*}
In a finite alphabet these follow by taking formal logarithms of
\(\prod_i(1-x_it)^{-1}\) and \(\prod_i(1+x_it)\), respectively.
Thus \(\omega(h_k)=e_k\), also with the zero convention in
negative degree.

On a permutation of cycle type \(\nu\), the sign is
\(\prod_{j\in\nu}(-1)^{j-1}\), counting parts with multiplicity.
The characteristic map of \ref{res:C05} therefore intertwines
tensoring with the sign representation and applying \(\omega\).
By \ref{res:A09} it follows that
\(\omega(s_\mu)=s_{\mu^{\mathsf t}}\). Apply \(\omega\) to the
ordinary Jacobi--Trudi formula of \ref{res:C04} for
\(\mu=\lambda^{\mathsf t}\). Use \(q\) rows in the determinant
before applying \(\omega\), and then specialize the stable
identity to the chosen finite alphabet. This gives the formula
for all alphabet sizes. See \cite[Theorem 4.5.1]{sagan}.
\end{proof}
